\documentclass[11pt]{article}
\usepackage[margin=27mm]{geometry}
\usepackage{amsmath,amssymb,amsthm,mathtools,booktabs,microtype,hyperref}
\newtheorem{theorem}{Theorem}[section]
\newtheorem{lemma}[theorem]{Lemma}
\newtheorem{proposition}[theorem]{Proposition}
\newtheorem{corollary}[theorem]{Corollary}
\newcommand{\Q}{\mathbb Q}
\newcommand{\Z}{\mathbb Z}
\newcommand{\F}{\mathbb F}
\newcommand{\Nm}{\operatorname{Nm}}
\newcommand{\rad}{\operatorname{rad}}
\newcommand{\hhat}{\widehat h}
\title{A prepared elliptic divisibility detector\\and its five-coordinate encoding}
\author{\(\Z\Q\) Diophantine Project}
\date{20 November 2026}
\begin{document}
\maketitle
\begin{abstract}
A rational function of two points on \(y^2=x^3-2x\) has square reduced
denominator when the second point is an integral multiple of the first.
An odd primitive denominator divisor gives the converse for sufficiently
large reference indices, conditional on ordinary rational \(abc\).
The function admits a fixed class-seven preparation for every integral
quotient after the reference point is chosen. Together with the published
intrinsic-pole theorem and a finite-algebra encoder, this gives a
five-coordinate subrelation of elliptic index divisibility in which every
nonzero integral quotient occurs. It need not accept every preassigned
divisible pair. It does not identify an external rational input with that
quotient and therefore does not lower a defining-variable bound for
\(\Z\subset\Q\).
\end{abstract}

\paragraph{Provenance.}
The surrounding project incorporates the constructions and notes of Gabe
Conant with Astra~6.2. The detector and normalization below are subsequent
project work. The arithmetic realization used for the terminal pole is
precisely the published class-seven form of Burungale's Theorem~B; no
stronger square-class premise is added.

\section{An elementary detector}

Let \(E_a:y^2=x^3+ax\), where \(a\in\Z\setminus\{0\}\), and let \(H,U\)
be non-torsion rational points. Define
\begin{equation}\label{eq:b}
 q=x(U),\qquad g=1+q^2,\qquad
 N=y(H)+q^2g,\qquad b(H,U)=\frac{N}{g^2}.
\end{equation}
The denominator \(g\) has no rational zero. Formal parameters are
\(\xi=-x/y\), and valuations satisfy \(v_p(p)=1\).

\begin{lemma}\label{lem:multiple}
Suppose \(U=[m]H\) for a nonzero integer \(m\). At every odd good prime
\(p\nmid a-1\), the valuation \(v_p(b(H,U))\) is not negative odd.
At any prime, including a bad prime, if \(H\) is formal of depth
\(e_H=v_p(\xi(H))\ge1\), then \(U\) is formal of depth \(e_U\ge e_H\), and
\begin{equation}\label{eq:near-one}
                         v_p(b(H,U)-1)=4e_U.
\end{equation}
\end{lemma}
\begin{proof}
For a formal point \(Q\) on this integral short Weierstrass model,
\[
 v_p(x(Q))=-2e_Q,\qquad v_p(y(Q))=-3e_Q.
\]
The formal multiplication series is integral, so \(e_U\ge e_H\).
If \(U\) is formal, \(q^2g\) has valuation \(-8e_U\). This is smaller
than \(v_p(y(H))\), whether \(H\) is formal or nonformal. Thus \(b\)
is a unit. When \(H\) is formal, the exact identity
\[
                   b-1=\frac{y(H)-1-q^2}{g^2}
\]
has numerator valuation \(-4e_U\), proving \eqref{eq:near-one}.

Suppose instead that \(U\) is nonformal. Then \(H\) is nonformal, and
both points have integral affine coordinates. A unit \(g\) makes \(b\)
integral. If \(v_p(g)>0\), then \(q^2\equiv-1\pmod p\). The reduction
of \(U\) is outside \(E[2]\): a two-torsion abscissa has \(q=0\) or
\(q^2=-a\), giving respectively \(g=1\) or \(g=1-a\).
Since multiplication preserves \(E[2]\), the reduction of \(H\) is
also outside \(E[2]\). Consequently \(y(H)\), and then \(N\), is a
unit. Therefore \(v_p(b)=-2v_p(g)\) is even.
\end{proof}

\begin{lemma}\label{lem:reject}
At a good prime, suppose that \(H\) is formal of positive odd depth
\(e_H\), whereas \(U\) is nonformal. Then
\[
                v_p(b(H,U))=-3e_H-2v_p(1+x(U)^2)<0
\]
is odd.
\end{lemma}
\begin{proof}
The summand \(q^2g\) is integral and \(y(H)\) has valuation
\(-3e_H<0\), so the latter dominates \(N\). The denominator subtracts
an even nonnegative valuation.
\end{proof}

The asymmetry in \eqref{eq:b} is deliberate. For true multiples, the
normalization approaches one uniformly in \(m\) at a prepared finite
place. Thus a fixed scaling works even when \(v_2(m)\) is arbitrarily large.

\section{Ordinary \(abc\) and odd primitive divisors}

We use the classical curve and point
\begin{equation}\label{eq:curve}
 E:y^2=x^3-2x,\qquad P_0=(2,2),\qquad
 P=[2]P_0=(9/4,-21/8).
\end{equation}
The Mordell--Weil group is
\(E(\Q)=\Z P_0\oplus\langle(0,0)\rangle\), as recorded for
LMFDB~256.b1, Cremona~256b1 \cite{LMFDB}. Exact rank bounds and the
torsion order can also be checked with PARI/GP; the companion audit
distinguishes those checks from the recorded generator assertion.

The elementary two-descent map \(x\bmod\Q^{\times2}\) takes the values
\(2\) and \(-2\) on \(P_0\) and \((0,0)\). Its image therefore has order
four, and its kernel is
\begin{equation}\label{eq:chart}
 2E(\Q)=\Z P,\qquad
 x=r^2,\quad y=r\eta,\quad \eta^2=r^4-2
\end{equation}
on the nonzero affine chart. The torsion point is excluded. The point
ordinate is nonzero on every rational solution of this chart.

We normalize the canonical height by
\(h(x(Q))=2\hhat(Q)+O(1)\).
The eventual odd-primitive-divisor assertion below is the denominator
case of Cornelissen--Zahidi, Proposition~3.10(iii) \cite{CZ}, which cites
Silverman, Lemma~13 \cite{Silverman}. Those results require ordinary
rational \(abc\), not an additional inert-prime conjecture. For this
specific curve one can make the \(abc\) step transparent as follows.

\begin{proposition}\label{prop:odd}
Assume ordinary rational \(abc\). Write
\[
 x([h]P)=(A_h/B_h)^2,\qquad
 \gcd(A_h,B_h)=1,\quad B_h>0.
\]
For every sufficiently large positive \(h\), the product \(Q_h\) of
primitive prime divisors of \(B_h\) occurring with exponent exactly one
satisfies
\begin{equation}\label{eq:primitive-mass}
                       \log Q_h\ge \frac15 h^2\hhat(P).
\end{equation}
In particular, one may require an odd primitive divisor outside any
prescribed finite set of primes.
\end{proposition}
\begin{proof}
The chart \eqref{eq:chart} supplies an integer \(C_h\) with
\[
 C_h^2+2B_h^4=A_h^4.
\]
If \(B_h\) were odd, reduction modulo \(16\) would make \(C_h^2\)
equal to \(14\) or \(15\), which is impossible. Thus \(B_h\) is even,
and \(A_h,C_h\) are odd. The three terms in this equation are coprime.
Also \(|C_h|\le A_h^2\) and \(|A_h|>B_h\).
Apply rational \(abc\) with \(\epsilon=1/100\):
\[
 |A_h|^4\ll \rad(2A_hB_hC_h)^{101/100}
          \le (2|A_h|^3\rad(B_h))^{101/100}.
\]
Writing \(L_h=\log|A_h|\), this gives
\begin{equation}\label{eq:radical}
                    \log\rad(B_h)\ge \frac{97}{101}L_h-O(1).
\end{equation}
Let \(V_h\) be the part of \(B_h\) supported on primes whose exponent
in \(B_h\) is at least two. Since
\(\rad(B_h)\le B_h/\sqrt{V_h}\), we obtain
\begin{equation}\label{eq:powerful}
                     \log V_h\le\frac8{101}L_h+O(1).
\end{equation}

Let \(D_h\) be the full nonprimitive part of \(B_h\). At a good odd
prime whose rank of apparition is \(r\mid h\), the formal-group identity
is
\[
        v_p(B_h)=v_p(B_r)+v_p(h/r).
\]
For a nonprimitive prime choose a prime divisor \(q\mid h/r\).
The term \(B_{h/q}\) has the same valuation, except that choosing
\(q=p\) can lose one. Consequently, up to the contribution of the
fixed bad primes,
\[
                    D_h\mid h\prod_{q\mid h,\ q\ {\rm prime}}B_{h/q}.
\]
The fixed-prime contribution is \(O(\log h)\) by the formal group
after entering its fixed open subgroup. Canonical-height bounds give
\[
 \log D_h\le h^2\hhat(P)\sum_{q\mid h}q^{-2}+O((\log h)^2)
          <\frac23 h^2\hhat(P)+O((\log h)^2).
\]
Here the sum is bounded by \(\zeta(2)-1<2/3\); no estimate uniform
in the curve is needed.

The classical Siegel approximation estimate for multiples of a fixed
non-torsion point gives
\[
        \log B_h\sim L_h\sim h^2\hhat(P).
\]
This is the same denominator-height input used in the proof of
\cite[Proposition~3.10]{CZ}. Removing the nonprimitive part and the
part \eqref{eq:powerful} leaves
\[
 \log Q_h\ge
 \left(1-\frac23-\frac8{101}-o(1)\right)h^2\hhat(P).
\]
The constant \(1-2/3-8/101=77/303\) exceeds \(1/5\), proving
\eqref{eq:primitive-mass}. A fixed finite set contributes at most
the sum of the logarithms of its primes to \(Q_h\), so it can be avoided.
\end{proof}

\begin{corollary}\label{cor:division}
Fix a finite set \(S\) containing the bad primes. There is an
ineffective \(h_0\) such that, for nonzero \(h,j\in\Z\) with
\(|h|>h_0\),
\[
 H=[h]P,\quad U=[j]P,\quad
 b(H,U)\text{ has no negative odd valuation outside }S
 \quad\Longrightarrow\quad h\mid j.
\]
\end{corollary}
\begin{proof}
Choose a primitive exponent-one prime \(p\notin S\) of \(B_{|h|}\).
At a good prime its rank of apparition is the reduction order
\(|h|\) of \(P\). Thus \(H\) is formal of odd depth, and if
\(h\nmid j\), then \(U\) is nonformal. Lemma~\ref{lem:reject}
contradicts the assumed denominator condition. Signs do not affect
the denominator or reduction order.
\end{proof}

\section{A fixed class-seven preparation}

Put
\begin{equation}\label{eq:d}
               d(H,U)=\frac{112}{841}b(H,U)
                     =7(4/29)^2b(H,U),\qquad
               S=\{2,3,7,29\}.
\end{equation}
For \(a=-2\), the extra excluded prime in Lemma~\ref{lem:multiple}
is \(3\). The scaling does not change valuations outside \(S\).

\begin{proposition}\label{prop:prepare}
For every nonzero \(m\in\Z\), there are arbitrarily large positive
\(h\) such that \(H=[h]P\), \(U=[m]H\), and
\begin{equation}\label{eq:prepared}
 0<d<1,\qquad d\in7\Q_2^{\times2},\qquad v_2(d)=4,\qquad
 d\text{ has square reduced denominator}.
\end{equation}
The reference can avoid any fixed finite set and all nonidentical
point guards used below.
\end{proposition}
\begin{proof}
Require \(H\) to be formal at all primes in \(S\). This is one
divisibility condition on \(h\), because the relevant local
quotients are finite. Equation~\eqref{eq:near-one} yields
\(b\in1+p^4\Z_p\) at these primes, uniformly for the integral
multiplier \(m\). In particular \(b\) is a two-adic square.
The valuations of \(d\) at \(2,3,7,29\) are respectively
\(4,0,1,-2\). At every other prime, Lemma~\ref{lem:multiple}
shows that its negative valuations are even.

Fix a compact real arc \(I\) around \(-P\) in the identity
component, small enough that \(0<y(H)<3\) on \(I\). The formula
\[
 b(H,U)=\frac{y(H)}{(1+x(U)^2)^2}
             +\frac{x(U)^2}{1+x(U)^2}
\]
gives \(0<b<4\) there, independently of \(m\). Thus
\(0<d<448/841<1\).
Positive multiples in the fixed divisibility progression are dense
in the real identity component: this is the elementary density of
a non-torsion element in the circle, and remains true after taking
any nonzero fixed multiple. Such \(h\) therefore enter \(I\)
arbitrarily late.

On the curve \(U=[m]H\), the signed numerators
\[
             x([m]H)^2(1+x([m]H)^2)\pm y(H)
\]
are nonzero rational functions. At \(H=O\), their two summands
have pole orders eight and three. Each therefore has only finitely
many zeros. The point-map denominators and every other already
verified nonzero rational guard can likewise be avoided without
losing the choice.
\end{proof}

This is completeness for every quotient after choosing its reference.
It is not a characterization of all preassigned divisible pairs:
the fixed scaling need not prepare every pair \((H,[m]H)\).

\section{The finite algebra and five coordinates}

The intrinsic pole theorem \cite{Intrinsic,Burungale} uses
\begin{equation}\label{eq:intrinsic}
 Y^2=P_d(v,s)
   :=d(d+v^2)^2s^4-d(d-v^2)^2(1-d)^2.
\end{equation}
Its rational solutions exclude negative odd valuations of \(d\)
at every prime, including all zero-coordinate cases. For inputs
\eqref{eq:prepared}, the published class-seven Theorem~B supplies
rational solutions with Zariski-dense projection to \((v,s)\).
The complete verification of its square-class and local hypotheses
is the one in \cite{Intrinsic}; neither a beta-one extension nor
a rational point on a prescribed fiber is inferred from density.

Use the two charts
\[
 H=(r^2,r\eta),\quad U=(t^2,t\theta),\qquad
 \eta^2=r^4-2,\quad \theta^2=t^4-2.
\]
Set
\begin{equation}\label{eq:AG}
 G=29(1+t^4),\qquad A=112(r\eta+t^4+t^8),\qquad d=A/G^2.
\end{equation}
The invertible change \(Z=G^5Y\) clears \eqref{eq:intrinsic} to
\begin{equation}\label{eq:terminal}
\begin{split}
 Z^2={}&A G^4(A+G^2v^2)^2s^4\\
       &-A(A-G^2v^2)^2(G^2-A)^2=:R.
\end{split}
\end{equation}
The two point equations and this monic quadratic equation define
a finite algebra of rank eight over \(\Q[r,t,v,s]\).
The terminal radicand has total degree \(48\) and \(162\) monomials
before reducing powers of \(\eta\). Its weighted degree is also
\(48\) when \(\eta,\theta\) have weight two.

For soundness exclude the finite index set \(0<|h|\le h_0\) by
\[
                         \prod_{k=1}^{h_0}(r^2-x([k]P))\ne0.
\]
This product, the point radicands, and the two signed values
of \(A\) enter the scalar guard. The signed-\(A\) zeros are
avoided in Proposition~\ref{prop:prepare}. For the resulting
fixed point algebra, the two signed values of the terminal
radicand are nonzero polynomials in \(v,s\): their leading
\(s^4\) coefficient is nonzero because the signed \(d\) is
nonzero. Intrinsic-pole density therefore avoids their product.
This makes the full algebra etale without adding a free coordinate.

For completeness of the compression argument, write
\(\alpha_j=\eta+j\theta+j^2Z\), \(1\le j\le57\).
Any pair of distinct geometric characters excludes at most two
values of \(j\), since its difference is a nonzero polynomial
of degree at most two. There are \(\binom82=28\) pairs.
Hence one member separates all eight characters.
Let \(\Delta_j\) be its power-basis determinant, \(M_j\) its
multiplication matrix, and \(c_j\) the product of \(\Delta_j\)
with the scalar guards. After clearing their fixed rational
coefficient denominators, define
\[
                         E_j(T)=\det(c_j(TI-M_j)-I).
\]
If \(c_j=0\), then \(E_j=1\). If \(E_j(T)=0\), then
\(c_j\ne0\) and \(T-1/c_j\) is a rational eigenvalue of a
primitive element, giving a rational character of the algebra.
Conversely a rational character gives the root
\(T=\chi(\alpha_j)+1/c_j\). The product of the \(57\) factors
therefore encodes the guarded source with one additional
rational variable.

\begin{theorem}\label{thm:subrelation}
Assume ordinary rational \(abc\) and the published class-seven
intrinsic-pole realization. The preceding finite specification
uses five total rational coordinates \(r,t,v,s,T\).
Every solution represents nonzero \(H=[h]P,U=[j]P\) with
\(h\mid j\). Every nonzero integral quotient \(m\) occurs in a
solution with \(U=[m]H\) for a suitable reference \(H\).
The specification need not accept every fixed divisible pair.
\end{theorem}
\begin{proof}
A rational character gives \eqref{eq:intrinsic}, so \(d\) has
no negative odd valuation. Outside \(S\), the same holds for
\(b\). The finite index guard and Corollary~\ref{cor:division}
give \(h\mid j\).
For a prescribed nonzero \(m\), choose the reference by
Proposition~\ref{prop:prepare}, then use intrinsic-pole density
and the separating menu as just described.
\end{proof}

\section{The unresolved input-projection step}

Theorem~\ref{thm:subrelation} does not identify a free rational
number \(X\) with \(j/h\). Its five-coordinate count is a count
for a subrelation on elliptic point charts, not a count for a
polynomial with input projection \(\Z\).

The most direct additional comparison is
\(\xi(U)/\xi(H)-X\). Once \(U=[m]H\), its formal logarithmic
limit is \(m-X\), but a local divisibility conclusion leaves
an upper height term \(\log(1+|m|)\).
A contact height carrier involving \(U\) grows on the order
of \(m^2h^2\), whereas the primes supporting the one-sided
formal comparison come from the denominator of \(H\), of
logarithmic size on the order of \(h^2\). This cannot supply
uniform completeness for unbounded \(m\).
A carrier involving only \(H\) and \(X\) avoids that support
defect but leaves \(\log(1+|m|)\) uncontrolled.

An independent point \(A=[n]P\) and a further divisibility
condition \((h-1)\mid(j-n)\) yield only
\[
                              m\equiv n\pmod{h-1}.
\]
No sufficient bound on \(m-n\) has yet been established.
Density, local approximation, and the two divisibilities alone
do not turn this congruence into equality.
Consequently no nine- or ten-variable definition of the
ordinary integers is asserted here.

\section{Exact checks and their limits}

The companion \nolinkurl{check_eds_detector.py} is read-only by
default. It verifies the universal \(G^{10}\) clearing,
the identity for \(b-1\), the degree and coefficient expansion
of \eqref{eq:terminal}, \(288\) exact true-multiple fixtures,
primitive-odd rejection fixtures, and six fully prepared
inputs at \(H=[68]P_0\).
The optional \nolinkurl{--pari} mode checks exact rank bounds,
torsion order, and conductor using PARI/GP.
The finite computations do not prove \(abc\), the infinite
odd-primitive theorem, the intrinsic-pole realization, or
the unresolved input-projection step.

\begin{thebibliography}{9}
\bibitem{CZ}
G. Cornelissen and K. Zahidi,
\emph{Elliptic divisibility sequences and undecidable problems about
rational points}, J. reine angew. Math. 613 (2007), 1--33.
Version used: arXiv:math/0412473v3, Proposition~3.10 and its proof.
\url{https://arxiv.org/pdf/math/0412473}.
\bibitem{Silverman}
J. H. Silverman, \emph{Wieferich's criterion and the abc-conjecture},
J. Number Theory 30 (1988), 226--237, Lemma~13.
\url{https://doi.org/10.1016/0022-314X(88)90019-4}.
\bibitem{LMFDB}
The LMFDB Collaboration, \emph{Elliptic curve 256.b1}
(Cremona 256b1), consulted 20 November 2026.
\url{https://www.lmfdb.org/EllipticCurve/Q/256/b/1}.
\bibitem{Burungale}
A. Burungale, \emph{A Diophantine definition of \(\Z\) in \(\Q\)},
arXiv:2611.16130v1, Theorem~B.
\url{https://arxiv.org/abs/2611.16130}.
\bibitem{Intrinsic}
\(\Z\Q\) Diophantine Project, \emph{An intrinsic pole block with two
parameters}, 17 November 2026, v56.
\nolinkurl{papers/pole_parameters/intrinsic_pole_20261117_v56.tex}.
\end{thebibliography}
\end{document}
